Le flux de communication du système débute par la \textbf{collecte des données} à partir des moniteurs HealForce ou des machines anciennes (via OCR). Ces données sont ensuite \textbf{transmises au backend} via le réseau de surveillance interne. Le backend \textbf{analyse les données en temps réel} et les stocke dans une base de données MySQL. En cas de détection d’anomalies (par rapport aux seuils prédéfinis), le système \textbf{envoie des alertes} au personnel médical via e-mail, SMS ou notifications push (FCM). Le personnel médical peut alors \textbf{consulter les alertes et les données} via l’application web ou mobile, et \textbf{prendre les mesures nécessaires} pour répondre aux situations critiques.

Ce flux, illustré dans la figure~\ref{fig:communication_flow}, assure une communication fluide et efficace entre les appareils médicaux, le système central et le personnel médical.
